\begin{textsample}{POS Dim 1 – pr – Score 64.00 – pr/pr\_em\_88.txt}  \label{ex:f1_pos_023}
5.
AVALIAÇÃO EM CIÊNCIAS DA NATUREZA E SUAS TECNOLOGIAS A globalização trouxe à sociedade \textbf{moderna} o desenvolvimento econômico, tecnológico e a vivência multicultural, colocando educadores, estudantes e a escola à frente de \textbf{um} mundo com novas exigências.
Nesse sentido, as relações pedagógicas estabelecidas na rotina diária escolar \textbf{que} \textbf{permeiam} o trabalho pedagógico como o planejamento, o currículo e a avaliação, apresentam diversos questionamentos e indagações sobre os novos desafios educacionais.
Toda e qualquer ação educadora, sistêmica e coerente, obrigatoriamente se depara com reflexões e discussões sobre o processo avaliativo.
O ato de avaliar a apropriação de conteúdos escolares só faz sentido se compreendido e realizado segundo o desenvolvimento de conceitos essenciais para \textbf{que}, assim, a avaliação se torne \textbf{uma} atribuição de qualidade pautada em dados relevantes de aprendizagem para os estudantes, auxiliando -os em tomadas de decisões ( Luckesi, 2011 ).
Desse modo, a avaliação escolar deve ser orientada para todo o processo de \textbf{ensino-aprendizagem} e não somente para os resultados ( Zabala e Arnau, 2010 ), cumprindo, assim, com a função e/ou finalidade integrada ao Projeto Político-Pedagógico ( PPP ) da escola, alinhada aos objetivos da formação, tendo a aprendizagem como processo educativo e pedagógico.
Outra \textbf{abordagem} a ser considerada é \textbf{que} a avaliação deve estar articulada ao currículo, ao ensino e à formação integral do \textbf{sujeito}, a fim de torná -lo protagonista, \textbf{um} estudante \textbf{pesquisador} e construtor do seu conhecimento.
Dessa maneira, o estudante poderá ser capaz de refletir e avaliar sobre a concepção de realidade em \textbf{uma} \textbf{totalidade} integrada, socialmente e \textbf{historicamente} situada no ambiente em \textbf{que} desenvolve ações participativas de forma crítica e sensível às questões ambientais e socioculturais na comunidade em \textbf{que} \textbf{vive}.
É importante \textbf{que} se considere, ao longo do processo de \textbf{ensino-aprendizagem}, \textbf{uma} avaliação formativa \textbf{que} traga ao estudante e ao docente maneiras de refletir sobre o aprendizado e as possíveis lacunas ou falhas durante determinado período.
Por meio da avaliação formativa, o objetivo do processo deixa de ser \textbf{uma} \textbf{mera} verificação dos conteúdos e passa a comandar a gestão da sala de aula, \textbf{uma} vez \textbf{que} devolve ao professor \textbf{uma} síntese do \textbf{que} foi aprendido e melhor aproveitado pelos estudantes. \textbf{Um} caráter individualizado de análise também nasce nessa perspectiva, pois a avaliação deve se adaptar ao processo na qual está inserida, levando em consideração a bagagem de conhecimento trazida pelos estudantes, contrapondo com o conhecimento escolar adquirido.
Ainda sob esse ponto de vista, a avaliação deve ser contínua, presente em diversos momentos ao longo do período letivo, e não pontual ao final do processo ( Barreira ; Boavida ; Araújo, 2006 ).
É necessário tanto para o educador quanto para o estudante compreender \textbf{que} o processo avaliativo curricular acarreta \textbf{uma} avaliação formativa.
O ato de avaliar deve se \textbf{configurar} como parte integrante do processo ensino-aprendizagem, sendo \textbf{imprescindível} no \textbf{direcionamento} das ações pedagógicas.
É indispensável, como \textbf{método} educativo, a intenção de melhorá -lo continuamente, principalmente quando se entende \textbf{que} estão intrinsecamente articulados ao desenvolvimento formativo educacional em todos os aspectos pedagógicos, tais como avaliação, aprendizagem e construção do conhecimento escolar.
Outro aspecto importante a ser considerado é \textbf{que}, para o ensino de Ciências da Natureza e suas Tecnologias ( CNT ), não podemos falar em mudanças pedagógicas na concepção de ensino para a área se o educador e a escola não tiverem a sensibilidade para perceber tais alterações.
Se a avaliação persistir em ser tradicional, seletiva, classificatória, somativa como única maneira de avaliar, concebida como juízo de valor do docente ou \textbf{um} valor atribuído como nota, não traduz a subjetividade da real apropriação dos conteúdos escolares pelos estudantes ( Luckesi, 2011 ).
O docente precisa se entender no processo avaliativo como balizador e orientador das potencialidades do processo de \textbf{ensino-aprendizagem}.
Segundo Sacristán e Gómez ( 1998 ), necessária a compreensão de \textbf{que} avaliar não é \textbf{uma} prática para comprovar rendimento ou peculiaridades dos estudantes, mas sim o fechamento de \textbf{um} ciclo avaliativo planejado, desenvolvido e analisado.
A avaliação é \textbf{um} instrumento para planejar e pensar a prática pedagógica.
Cabe \textbf{salientar} \textbf{que} a avaliação formativa é aquela \textbf{que} ocorre no decorrer das etapas educacionais com o objetivo de reorientar a prática docente perante o acompanhamento do desenvolvimento do estudante ( Sacristán e Gómez, 1998 ).
Segundo Zabala : > A partir de \textbf{uma} opção \textbf{que} contempla como finalidade fundamental do ensino a formação integral da pessoa, e conforme \textbf{uma} concepção construtivista, a avaliação sempre deve ser formativa, de maneira \textbf{que} o processo avaliador, independentemente de seu objeto de estudo, tem \textbf{que} observar as diferentes fases de \textbf{uma} intervenção \textbf{que} deverá ser estratégica ( Zabala, 1998, p. 201 ).
O processo avaliativo no decorrer da etapa do Ensino Médio emerge dos pressupostos legais das Diretrizes Curriculares Nacionais para o Ensino Médio ( DCNEM ), da Lei de Diretrizes e Bases da Educação Nacional ( LDBEN ), n. 9.394/96, e da Base Nacional Comum Curricular ( BNCC ), haja vista \textbf{que} os estudantes se apropriam do conhecimento escolar de diferentes formas, sendo cada indivíduo único, sendo \textbf{preciso} respeitar a individualidade e o ritmo de aprendizagem de cada \textbf{um}.
As propostas avaliativas ofertadas por políticas públicas \textbf{contemporâneas} visam garantir a qualidade do trabalho escolar, porém essa concepção apresenta múltiplos significados, quer seja no âmbito pedagógico, quer seja no âmbito social e político, todavia não podem ser desvinculadas da \textbf{totalidade} do contexto histórico e das relações sociais.
De acordo com as Diretrizes Curriculares Nacionais Gerais da Educação Básica, a avaliação educacional, por meio da \textbf{ótica} \textbf{teórica}, é apresentada em três dimensões básicas : a ) avaliação da aprendizagem ; b ) avaliação institucional interna e externa ; c ) avaliação de redes de Educação Básica ( Brasil, 2013 ).
A avaliação da aprendizagem é constituída por conjuntos de habilidades, conhecimentos, princípios, valores direcionados aos \textbf{sujeitos} associados de maneira integrada às concepções da Educação Básica. \textbf{Necessita} ser concebida como processo avaliativo integrado e articulado a \textbf{uma} \textbf{abordagem} avaliativa formativa, como prática pedagógica investigativa, \textbf{dialógica}, participativa, \textbf{que} deve apresentar \textbf{um} diagnóstico preliminar e ser entendida como parte do processo avaliativo educacional. \textbf{Salienta} -se \textbf{que} neste Referencial Curricular para o Ensino Médio, direcionado à área de conhecimento Ciências da Natureza e suas Tecnologias, prioriza -se a avaliação da aprendizagem significativa.
Pelizzari afirma \textbf{que} a aprendizagem “ é muito mais significativa à medida \textbf{que} o novo conteúdo é incorporado às estruturas de conhecimento do estudante e adquire significado para ele a partir da relação com seu conhecimento prévio ” ( Pelizzari \textbf{et} al., 2002, p. 38 ).
Diante do \textbf{exposto}, a literatura e os documentos \textbf{norteadores}, dentre eles a BNCC, apresentam como consenso a necessidade de proporcionar ao estudante o ensino por competências, sendo \textbf{que} as competências são expressas por conjuntos de habilidades, conhecimentos e atitudes \textbf{que} propõem ao estudante : sapiência na tomada de decisões, mobilização de recursos e ativação de hábitos educacionais em contextos de complexidade.
Ensinar por competências é ter, antes de tudo, \textbf{uma} relação intrínseca com o conteúdo apresentado em sala de aula, com o objetivo de desenvolver nos estudantes as habilidades, a partir do currículo previamente proposto pela escola e, consequentemente, as competências almejadas.
É \textbf{preciso} inovar nas práticas e no planejamento, promovendo o ensino por meio da \textbf{contextualização}, da resolução de situações-problema, do processo investigativo, além da integração do indivíduo com o conhecimento científico associada às trajetórias de vivências desse estudante.
Avaliar no nível das competências é \textbf{uma} \textbf{tarefa} complexa, visto \textbf{que} provoca \textbf{uma} análise de situações-problema, partindo de contextos reais para o aperfeiçoamento das habilidades vinculadas à competência \textbf{que} se deseja quantificar ( Zabala e Arnau, 2010 ).
Desse modo, as habilidades desenvolvidas avaliam situações reais \textbf{que} colocam o estudante o mais próximo possível de circunstâncias \textbf{que} ele poderá enfrentar fora do ambiente escolar.
Tendo em vista \textbf{que} os exames externos ainda estejam presos a \textbf{provas} escritas, há dificuldade em analisar as habilidades como parte de todo processo de \textbf{ensino-aprendizagem} e, consequentemente, as competências vinculadas, fora do escrever e do calcular.
A instituição escolar deve oferecer \textbf{um} amplo leque de instrumentos \textbf{que} possibilitem avaliar a oralidade, a capacidade em tomar decisões, de enfrentar \textbf{crises}, de levantar hipóteses, dentre outras habilidades \textbf{que} são desenvolvidas durante o período escolar.
Ao realizar a avaliação, o conhecimento escolar adquirido ao longo do processo de \textbf{ensino-aprendizagem} deve ser priorizado.
As \textbf{provas} e testes limitados a questões de múltipla escolha e discursivas \textbf{que} direcionam as respostas do estudante de acordo com a esperada pelo docente não podem ser mantidas como únicos instrumentos de avaliação.
Avaliar se o estudante desenvolveu determinada competência ou habilidade vai além da \textbf{prova} escrita.
Primeiramente, é \textbf{preciso} conhecer o \textbf{que} o estudante já sabe, e a partir desse ponto, criar mecanismos \textbf{que} explorem as habilidades desenvolvidas e a \textbf{aplicabilidade} do conhecimento \textbf{que} deve ser repensado, não apenas no momento da avaliação, mas durante todo o processo de \textbf{ensino-aprendizagem}.
Ao pensar o ensino das Ciências Naturais sob o olhar da aprendizagem significativa, a avaliação precisa estar presente desde o início do processo, quando é necessário \textbf{um} diagnóstico sobre o conhecimento do estudante, sua condição social, seus anseios e objetivos.
É durante o processo, através da avaliação formativa, \textbf{que} se busca \textbf{uma} adequação da metodologia para cada indivíduo ou grupo, identificando ao longo do caminho os estudantes em diferentes níveis de aprendizagem, não com o sentido de classificar e selecionar, mas de adequar o plano de trabalho docente de acordo com cada grupo de estudantes.
Assim, a avaliação passa a diagnosticar a situação de ensino em \textbf{que} se encontra cada estudante, retornando atribuições conforme a situação de cada \textbf{um} e contribuindo para a modificação e melhoria nas ferramentas educacionais utilizadas pelo docente.
E por fim, \textbf{uma} avaliação ao final do processo, avaliando o domínio do estudante sobre objetivos previamente determinados e combinados, indicando ao docente o \textbf{que} foi ou não aprendido durante o período letivo ( Villattorre ; Higa ; Tychanowicz, 2008 ).
Com essa característica apresentada, a avaliação faz parte da aprendizagem significativa, pois busca no estudante conteúdos com os quais ele já está familiarizado, ideias âncoras a partir das quais construir novos conhecimentos, adaptando e aumentando a estrutura cognitiva.
Quando o estudante percebe \textbf{uma} relação entre o \textbf{que} ele aprende na escola e o \textbf{que} ele \textbf{vive} em seu cotidiano, o conhecimento ganha raízes em sua estrutura cognitiva facilitando assim o seu aprendizado.
As avaliações devem trazer, portanto, elementos \textbf{que} \textbf{lembrem} a utilização prática dos conceitos aprendidos, contextualizando e explorando o cotidiano.
A intencionalidade deste Referencial Curricular para o Ensino Médio é propor \textbf{que} a área de conhecimento Ciências da Natureza e suas Tecnologias aproprie -se de \textbf{uma} avaliação \textbf{que} esteja centrada no desenvolvimento do letramento científico, com base no ensino por investigação articulado à \textbf{contextualização}, \textbf{que} pode se dar por meio do enfoque CTS ( Ciência - Tecnologia - Sociedade ), buscando \textbf{uma} avaliação integrada à \textbf{abordagem} formativa e emancipadora, fundamentada em \textbf{uma} aprendizagem \textbf{que} seja significativa ao estudante.
Espera -se, assim, \textbf{que} o processo avaliativo não contemple somente resultados, mas \textbf{que} se torne \textbf{um} potencial instrumento de aprendizagem com o intuito de ser mais \textbf{dialógico}, participativo, democrático e \textbf{que} desconstrua a avaliação tradicional pautada, por vezes, em \textbf{um} único instrumento avaliativo.
A BNCC propõe \textbf{que} o ensino de Ciências da Natureza e suas Tecnologias deve ir além do aprendizado de seus conteúdos essenciais, sendo assim é \textbf{preciso} discutir o papel do conhecimento científico e tecnológico na organização social, nas questões ambientais, na saúde humana e na formação cultural ou seja, analisar as relações entre ciência, tecnologia, sociedade e ambiente ( Brasil, 2018 ).
Dessa forma, a aprendizagem deve valorizar a aplicação dos conhecimentos na vida individual, favorecendo o protagonismo dos estudantes no enfrentamento de questões cotidianas.
Para tanto, é \textbf{preciso} vincular as CNT a \textbf{tópicos} socioambientais e sociotecnológicos como temáticas para resolução de situações-problema e tomada de decisões \textbf{que} articulem os conhecimentos escolares e as vivências diárias do estudante, tendo como ponto \textbf{central} os conteúdos essenciais de cada componente curricular desta área de conhecimento, as competências específicas da área e as habilidades a serem desenvolvidas.
Nesse contexto, a avaliação possibilita a identificação de obstáculos, a determinação de objetivos e o planejamento de ações a serem tomadas ao longo do processo de \textbf{ensino-aprendizagem}, beneficiando as habilidades desenvolvidas pelos estudantes.
Sob a concepção da aprendizagem significativa, o estudante é colocado como \textbf{sujeito} ativo no processo de aprendizagem, apto a se apropriar de novos conhecimentos, de racionalizar e de criar interação entre os conhecimentos previamente adquiridos em suas vivências com o novo conhecimento, ou seja, as informações recentes se tornam \textbf{expressivas} e significativas em \textbf{uma} organização cognitiva preexistente do estudante, possibilitando a construção de \textbf{uma} base de \textbf{fundamentos} específicos, sendo incentivado a pensar, explorar significados, buscar diferenças e paridades a fim de alcançar a aprendizagem.
Dessa maneira, é \textbf{preciso} \textbf{que} o docente compreenda \textbf{que}, para se atingir \textbf{uma} aprendizagem significativa, é necessário \textbf{que} haja \textbf{um} ensino centrado no estudante, \textbf{que} apresente mudanças na prática pedagógica do educador, como o conhecimento do público alvo, \textbf{que} contribui para transformá -lo em \textbf{um} protagonista, em \textbf{um} \textbf{pesquisador} do conhecimento científico apto a perceber a importância e a \textbf{correlação} entre as práticas escolares e as sabedorias apropriadas em sua rotina diária.
Para \textbf{que} a aprendizagem significativa ocorra é necessário \textbf{que} sejam feitas, com frequência, avaliações diagnósticas \textbf{que} apontem quais são as potencialidades e fragilidades no conhecimento do estudante.
Para as CNT é possível a utilização de \textbf{encaminhamentos} \textbf{metodológicos} \textbf{que} envolvam mapas \textbf{conceituais}, resolução de situações-problema, ensino por investigação, \textbf{contextualização} dos conceitos científicos articulados à \textbf{interdisciplinaridade}, \textbf{abordagem} por meio do enfoque CTS, da História da Ciência, dos três momentos pedagógicos, \textbf{que} têm como finalidade tornar os conhecimentos científicos significativos e não isolados, além de incorporá -los aos conhecimentos escolares e aos conceitos prévios.
A dissociação do conteúdo escolar do contexto sociocultural, econômico, político, ambiental e científico, bem como a ausência de significado dos conhecimentos e, consequentemente, a falta de sentido dos conceitos ensinados na escola, desvinculados dos interesses dos estudantes, fazem parte de \textbf{um} discurso \textbf{que} aponta \textbf{uma} \textbf{crise} no Ensino de Ciências ( Neto, 1995 ; Fourez, 2003 ; Delizoicov e Angotti, 2002 ; Pozo e Crespo, 2009 ).
Assim, emerge o campo do conhecimento Didática das Ciências.
Nos \textbf{dizeres} de Astolfi e Develay, a definição de Didática das Ciências surge do entendimento de Bachelard, Piaget e Bruner, \textbf{uma} vez \textbf{que} o interesse desse campo de conhecimento é associar a forma de ensinar com os processos de aprendizagem, objetivando o aprofundamento dos conhecimentos e tornando -os mais significativos aos estudantes ( Astolfi ; Develay, 1995 ), tendo em vista \textbf{que} toda aquisição de conceitos interfere nas experiências preexistentes ( Cachapuz \textbf{et} al., 2001 ).
Ainda para esses \textbf{autores}, dessa forma possibilita -se ao docente \textbf{uma} explicação eficaz e \textbf{funcional} de fenômenos naturais e conceitos científicos de maneira a oferecer significância às concepções já existentes no sistema cognitivo dos estudantes.
Assim é \textbf{que} o campo da Didática das Ciências está em consonância com os objetivos da aprendizagem significativa.
No entanto, cabe à escola e ao educador a apropriação de novas metodologias e mudanças no modo de conceber o ensino em ciências, oportunizando \textbf{uma} aprendizagem mais prazerosa e efetiva.
Tal perspectiva alinha -se à proposta de \textbf{uma} escola democrática e inclusiva, \textbf{que} prioriza o conhecimento científico percebendo \textbf{que} currículo e avaliação são \textbf{categorias} indissolúveis, pois encontram -se em \textbf{uma} interação constante e dinâmica.
O percurso do desenvolvimento curricular acarreta ao educador diversos aspectos destinados a avaliar as aprendizagens dos estudantes, porém precisa -se compreender a trajetória de vivências desses jovens e a concepção de educação, currículo, ensino e avaliação oriunda do sistema educacional, para \textbf{que}, assim, construam -se expectativas sobre quais os conhecimentos serão avaliados, valorizando o \textbf{que} de fato o estudante está adquirindo dentre os conteúdos ministrados.
Nas Ciências da Natureza e suas Tecnologias é necessário \textbf{que} o docente conheça e organize as habilidades a serem desenvolvidas, partindo da seleção dos conteúdos curriculares essenciais, vinculadas aos \textbf{encaminhamentos} \textbf{metodológicos}.
Em seguida, define -se quais instrumentos avaliativos da aprendizagem serão empregados ao conhecimento do estudante, tornando -o parte integrante do processo, evitando -se, assim, enunciados dúbios, ambíguos, questões meramente copiadas de bancos de dados sem as devidas adequações ao currículo real ministrado em sala de aula, compreendendo o \textbf{que} se quer cobrar e como cobrar.
Caso contrário, as discussões sobre avaliação e currículo se tornam desnecessárias.
O ensino por meio da investigação científica facilita a \textbf{problematização} de situações cotidianas para elucidar \textbf{um} dado fenômeno, o \textbf{que} \textbf{desperta} nos estudantes a necessidade de busca por informações para analisar o problema a ser resolvido quando suas concepções prévias não forem suficientes.
Ao elaborar seus objetivos, os docentes devem relacioná -los com a intenção de desenvolver as habilidades definidas para cada componente curricular da área.
Nesse sentido, a prática da avaliação no processo educativo de diagnosticar, perpassa o acompanhamento do processo de \textbf{ensino-aprendizagem}, \textbf{configurando} -se como \textbf{um} importante instrumento de investigação da prática pedagógica, fornecendo informações necessárias para a mediação entre as intervenções do docente e o estudante.
A avaliação também auxilia a ação do docente, provocando indagações sobre como ele estimula, \textbf{instiga} e motiva o estudante, bem como proporciona reflexões relacionadas ao diálogo, troca de ideias e apoio para tomada de decisões durante o processo educativo.
Por sua vez, a essência investigativa na área das Ciências da Natureza e suas Tecnologias favorece a relação entre a inserção dos conhecimentos escolares e as experiências de vida dos estudantes por meio de comparações e análises, empreendidas pelos mesmos, cujos resultados referem -se à ampliação dos significados sobre os fenômenos naturais e tecnológicos. \textbf{Uma} vez \textbf{que} a atividade investigativa esteja em foco no planejamento das aulas, o desenvolvimento do pensamento crítico dos estudantes é favorecido por meio de avaliações \textbf{que} privilegiem a observação sistemática, a elaboração de questionamentos, o levantamento de hipóteses e deduções \textbf{que} oportunizam o desenvolvimento da produção de relatos com o uso da linguagem científica.
Nesse contexto, avaliar dentro dessa área significa verificar a familiaridade do estudante com procedimentos, técnicas e equipamentos científicos, bem como seu \textbf{posicionamento} diante da \textbf{problematização}, suas atitudes apresentadas durante as atividades experimentais, sua apropriação da linguagem científica e das competências estabelecidas para os componentes curriculares da área, além da averiguação do progresso \textbf{conceitual} apresentado pelo estudante.
Dessa forma, o docente deve compreender \textbf{que} a finalidade de se avaliar é tão importante quanto conhecer os instrumentos avaliativos possíveis para \textbf{que} a avaliação seja diagnóstica e construtiva.
Entende -se por instrumentos avaliativos todos os recursos utilizados pelo docente \textbf{que} permitam a análise de dados referentes ao progresso da aprendizagem.
A diversificação desses instrumentos possibilita, tanto ao docente quanto ao estudante, oportunidades de refletir sobre seu avanço e de assegurar \textbf{que} o processo de \textbf{ensino-aprendizagem} esteja ocorrendo de forma \textbf{justa} e fidedigna às necessidades, oportunizando o aprofundamento e a retomada dos conteúdos sempre \textbf{que} necessário.
A apreciação construtiva sobre dados relevantes auxilia não só o docente e a escola na tomada de decisões sobre suas práticas, como também o estudante e a família na reflexão sobre seu progresso.
Quando a avaliação é parte de \textbf{uma} perspectiva de ensino \textbf{que} articula o conhecimento escolar com questões do cotidiano, colabora para \textbf{que} o estudante exerça sua cidadania de forma consciente perante a própria realidade e em relação aos acontecimentos globais.
Na sequência serão apresentadas algumas possibilidades avaliativas para cada componente curricular da área de Ciências da Natureza e suas Tecnologias.
A avaliação no componente curricular de Biologia está diretamente ligada à compreensão do fenômeno vida e suas relações com o ambiente, possibilitando \textbf{que} o estudante desenvolva a capacidade de observação dos fatos ao seu redor e, consequentemente, do conhecimento escolar adquirido.
Nesse sentido, a avaliação desenvolve no indivíduo a capacidade de compreender os conceitos básicos do componente curricular Biologia, pensando, adquirindo e avaliando informações de forma autônoma para aplicar seus conhecimentos no cotidiano e \textbf{despertar} seu interesse pelo mundo vivo.
Sendo assim, pode -se identificar \textbf{que} ocorreu o processo de alfabetização biológica, \textbf{que} segundo Krasilchik : > Admite -se \textbf{que} a formação Biológica contribua para \textbf{que} cada indivíduo seja capaz de compreender e aprofundar as explicações atualizadas de processos e de conceitos biológicos, a importância da ciência e da tecnologia na vida \textbf{moderna}, enfim, o interesse pelo mundo dos seres vivos ( Krasilchik, 2016, p. 12 ).
Krasilchik afirma também \textbf{que} existe \textbf{uma} necessidade de cautela no momento da escolha, da construção e da aplicação dos instrumentos de verificação do aprendizado e da análise dos resultados ( Krasilchik, 2016 ).
É primordial a preparação de instrumentos coerentes com os objetivos propostos pelo docente no seu planejamento curricular, como, por exemplo, a \textbf{abordagem} de temas de natureza polêmica, \textbf{que} vão do âmbito econômico, social, político, moral e até mesmo ao ético e religioso — como o uso dos transgênicos, a experimentação com animais, a reprodução humana assistida, a bioética, a biotecnologia, a problemática ambiental, a vacinação, o uso dos agrotóxicos, a clonagem, o aborto, as células-tronco, entre muitos outros —, possibilitando a realização de debates e outras atividades \textbf{que} permitam ao docente avaliar o desenvolvimento da consciência crítica e a condição argumentativa dos estudantes na tomada de decisões, sua formação ética e suas proposições quanto aos valores pessoais e sociais.
De acordo com Krasilchik e Marandino, os termos alfabetização e letramento científico possuem significados distintos, sendo \textbf{que} a expressão “ alfabetização científica engloba a ideia de letramento, entendida como a capacidade de ler, compreender e expressar opiniões sobre ciência e tecnologia [... ] ” ( Krasilchik ; Marandino, 2007, p. 18 ), permitindo assim \textbf{que} cada cidadão participe da cultura científica de forma individual e coletiva.
Segundo Scarpa e Silva, a Biologia possui princípios \textbf{que} a definem como \textbf{uma} ciência autônoma, impossibilitando -a de \textbf{uma} \textbf{abordagem} experimental de diversos temas dessa área de conhecimento, porém não inviabiliza a proposição de atividades de acordo com as propostas de \textbf{um} ensino por investigação para temas biológicos, contribuindo para \textbf{que} os estudantes desenvolvam diferentes habilidades presentes no fazer das ciências biológicas, franqueando -lhes, assim, a alfabetização científica ( Scarpa ; Silva, 2013 ).
No componente curricular Biologia, vários são os fatores \textbf{que} devem ser levados em conta pelo docente no processo avaliativo, como, por exemplo, a periodicidade das avaliações e os instrumentos avaliativos utilizados, o tempo de aprendizagem de cada estudante, a análise crítica e a criatividade ao selecionar os materiais, entre outros.
Diante disso, a busca por recursos didáticos \textbf{que} facilitem o processo avaliativo deve levar em consideração a importância \textbf{que} as TDIC ( Tecnologias Digitais de Informação e Comunicação ) exercem no processo de \textbf{ensino-aprendizagem}, através da utilização de modelos e jogos didáticos, plataformas on-line, blogs, sites, criação de vídeos e podcasts, mediação de aplicativos, gamificação, modelagem molecular, entre outros exemplos.
As atividades experimentais, estudos de caso, estudos do meio, seminários, debates, atividades lúdicas, leitura e interpretação de textos e imagens, entre outros, também possibilitam ao estudante o desenvolvimento de novos conceitos \textbf{que} buscam a construção do conhecimento científico.
Esses aspectos levam a \textbf{uma} reflexão na \textbf{abordagem} dos conteúdos escolares, nas escolhas das metodologias de ensino, dos critérios adotados nas avaliações e instrumentos avaliativos, \textbf{que} orientam os propósitos e as dimensões do \textbf{que} se avalia no componente curricular Biologia.
A avaliação no componente curricular de Física, assim como nos demais componentes da área de CNT, deve priorizar as habilidades específicas da área, concentrando nas \textbf{que} envolvem o componente.
Essas habilidades desenvolvidas não trazem apenas conhecimentos científicos, mas visam fatores sociais e ambientais.
Portanto, a avaliação deve trazer os conteúdos abordados e refletir sobre a prática social e ambiental \textbf{que} será \textbf{salientada} durante o processo.
Nesse sentido, ao avaliar, deve -se ter consciência em considerar os aspectos qualitativos sobre os quantitativos, como é descrito na LDB. \textbf{Uma} boa avaliação é aquela \textbf{que} contempla as habilidades \textbf{que} foram desenvolvidas ao longo do processo, possibilitando ao docente e ao próprio estudante \textbf{uma} reflexão sobre o processo de \textbf{ensino-aprendizagem}. \textbf{Uma} prática comum, \textbf{historicamente}, é fazer da avaliação \textbf{uma} \textbf{mera} verificação sobre aspectos do conhecimento \textbf{que} enfatizam \textbf{provas} escritas, resolução de exercícios, aplicação de fórmulas e verificação de \textbf{teorias}.
Cria -se, então, processos seletivos, classificatórios, com a intenção de treinar o estudante para vestibulares e exames de seleção, fazendo com \textbf{que} todo o currículo seja voltado para essa intenção.
Visando quebrar essa prática, é sugerido trabalhar com avaliações diagnósticas, identificando os conhecimentos prévios dos estudantes.
Fazendo o uso de mapas \textbf{conceituais}, debates, leitura crítica, por exemplo.
Os mapas \textbf{conceituais} como diagnóstico podem funcionar pedindo ao estudante \textbf{um} tema específico, porém abrangente, em \textbf{que} o estudante possa relacionar conceitos \textbf{que} ele já conhece de \textbf{uma} maneira \textbf{que} possibilite ao professor analisar qual é o nível de conhecimento e como esse se organiza na estrutura cognitiva do \textbf{sujeito}.
Já os debates e a leitura crítica podem trazer \textbf{uma} visão mais argumentativa do estudante sobre o tema, o \textbf{que} explora diversas habilidades presentes na área de CNT.
A aprendizagem significativa enfatiza \textbf{que} os novos conhecimentos adquiridos devem ter \textbf{uma} \textbf{correlação} com o novo conhecimento, a fim de consolidar o aprendizado na estrutura cognitiva já existente, e para \textbf{que} essa aprendizagem ocorra é \textbf{preciso} romper as barreiras entre o senso comum e o conhecimento científico.
Para isso são necessárias avaliações diagnósticas \textbf{que} explorem os conceitos a serem abordados e não a \textbf{simples} verificação de fórmulas e equações são fundamentais.
Ainda, torna -se indispensável \textbf{uma} avaliação contínua, com \textbf{uma} recuperação concomitante e, sempre \textbf{que} possível, fazendo a retomada dos conceitos \textbf{que} não foram compreendidos no primeiro momento. \textbf{Lembrando} \textbf{que} a avaliação como parte do processo também tem a intenção de moldar os \textbf{encaminhamentos} \textbf{metodológicos} aplicados.
Assim, o desenvolvimento das habilidades é melhor aproveitado tanto pelo estudante quanto pelo docente.
É necessário reforçar a ideia de \textbf{que}, para \textbf{que} a avaliação ocorra com caráter formativo é \textbf{preciso} \textbf{que} os \textbf{encaminhamentos} \textbf{metodológicos} utilizados nas aulas envolvam mais ferramentas do \textbf{que} simplesmente a explanação e transposição de conteúdos.
O docente deve ser o mediador da construção do conhecimento e fará do estudante protagonista nesse processo.
Assim, o estudante pode trazer seus conhecimentos prévios, suas dúvidas e curiosidades, \textbf{que} podem ser sanadas pelo conhecimento escolar adquirido, modificando e ampliando sua visão de mundo.
O desenvolvimento das habilidades durante as aulas não é só percebido por meio de avaliações pontuais realizadas ao final de determinados períodos, mas em todos os momentos.
Pensando nisso, as avaliações podem ser aplicadas concomitantemente ao desenvolvimento dessas habilidades.
Nesses momentos é possível fugir das tradicionais \textbf{provas} escritas, \textbf{que} \textbf{lembram} exames de vestibular, e \textbf{focar} em atividades \textbf{que} vão de encontro ao enfoque CTS, a \textbf{problematização}, a História e Filosofia da Ciência, as TDIC, entre outras estratégias \textbf{que} envolvam mais o estudante com os conteúdos apresentados e, por consequência, se \textbf{aproximam} das habilidades almejadas, fazendo com \textbf{que} o momento da avaliação também seja \textbf{um} momento de aprendizado.
Para Villattorre, Higa e Tychanowicz ( 2008 ), a avaliação quando tratada de maneira diagnóstica, formativa e somativa enriquece o processo de \textbf{ensino-aprendizagem}, pois indica o envolvimento dos estudantes com as habilidades e conteúdos trabalhados e o professor tem a possibilidade de analisar e rever quando necessário a didática utilizada.
Esses \textbf{autores} também enfatizam \textbf{que} a avaliação deve ser construída de maneira democrática, estabelecendo regras e critérios claros para todos.
Assim, o estudante saberá como será avaliado e quais os objetivos devem ser alcançados.
A Física é \textbf{vista} como \textbf{uma} ciência \textbf{que} estuda a natureza e seus fenômenos, e investiga a matéria, o espaço e o tempo.
O \textbf{que} estudamos são pedaços da realidade e, mesmo \textbf{que} não entendamos, eles continuam a existir.
São modelos \textbf{que} partem de ideias humanas, sem \textbf{uma} verdade absoluta ( Menezes, 2005 ).
A avaliação em Física deve ter o cuidado de não tratar questões como verdades absolutas, explorando a evolução dos conhecimentos, a divergência de ideias e a relação com os fatores sociais e ambientais.
O professor deve ter consciência \textbf{que} a avaliação nem sempre será objetiva, mas deve ganhar traços de subjetividade com instrumentos diversos \textbf{que} fogem do tradicional lápis e papel.
A avaliação no componente curricular de Química, como parte do processo educativo, quando imersa numa perspectiva de ensino \textbf{que} articula o conhecimento escolar com as vivências e questões do cotidiano, colabora para \textbf{que} o estudante exerça sua cidadania de forma consciente perante sua própria realidade e em relação aos acontecimentos globais e avanços tecnológicos.
Nesse sentido, a avaliação não deve se encerrar em momentos pontuais \textbf{que} \textbf{exigem} apenas \textbf{memorização} de fórmulas, símbolos, equações e resolução de exercícios mecânicos, tampouco se \textbf{restringir} a medir conhecimentos.
A avaliação é \textbf{um} meio contínuo de proporcionar suporte ao estudante no seu processo de apropriação dos conteúdos e no seu processo de constituição de si mesmo como \textbf{sujeito} existencial e como cidadão ( Luckesi, 2011 ).
Seguindo com os \textbf{dizeres} de Luckesi, ao utilizar instrumentos avaliativos como testes e \textbf{provas}, examina -se de forma classificatória e quantitativa a aprendizagem do \textbf{sujeito}, \textbf{baseando} -se somente na quantidade de \textbf{erros} e acertos.
Por sua vez, quando se avalia o estudante, presume -se diagnosticar e reorientar reforçando os pontos positivos e verificando os pontos negativos a fim de aprimorá -los ( Luckesi, 2011 ).
Na Educação Básica, o componente curricular Química surge no Ensino Médio com seus conhecimentos específicos sem \textbf{que} o educador faça \textbf{correlação} com conceitos escolares presentes na matriz curricular do Ensino Fundamental II.
No decorrer de sua vida escolar, os estudantes se deparam com conhecimentos \textbf{que} reproduzem instantaneamente e de maneira complexa os conceitos químicos, resultando, no final de sua formação básica, em \textbf{um} bloqueio de aprendizagem desses conhecimentos em sala de aula ( Kosminsky \& Giordan, 2002 ; Cunha \& Giordan, 2009 ; Siqueira, Silva e Júnior, 2011 ).
Faz -se necessária \textbf{uma} mudança dessa concepção para \textbf{que}, de acordo com esta proposta de Referencial Curricular para o Ensino Médio, haja \textbf{uma} continuidade e aprofundamento dos conhecimentos escolares do componente curricular Química, a partir dos conceitos \textbf{vistos} no componente Ciências.
Nessa perspectiva, a avaliação para este componente deve priorizar as metodologias e os instrumentos avaliativos \textbf{que} considerem os conhecimentos prévios provenientes da cultura do estudante, verificando em \textbf{que} medida este \textbf{consegue} articulá -los com a sua realidade social, transpondo -os para o conhecimento escolar. \textbf{Logo}, para este componente curricular, propõe -se a avaliação da aprendizagem significativa, diagnóstica e formativa como parte do processo educativo, considerando os \textbf{métodos} mais adequados para desenvolver habilidades de caráter atitudinais, investigativas, relacionadas à comunicação e expressão e argumentativas, pertinentes ao componente curricular Química.
Além disso, deve -se priorizar as unidades temáticas e os conteúdos essenciais preestabelecidos no currículo objetivando a análise da qualidade da aprendizagem dos conhecimentos apropriados pelos estudantes de maneira significativa, bem como o (re)planejamento das ações do educador em sala de aula.
Para tanto, os instrumentos de avaliação devem ser diversificados, considerando \textbf{que} cada \textbf{sujeito} aprende e expressa seus saberes de diferentes formas.
Ao oportunizar \textbf{que} os estudantes explorem seus conhecimentos por diversos meios, constroem -se cidadãos do processo democrático capazes, de maneira crítica, de tomada de decisões \textbf{que} envolvam situações-problema comuns à rotina diária articulando os conteúdos escolares e suas vivências.
Cabe ao corpo docente e pedagógico não limitar o processo avaliativo somente a \textbf{provas} escritas, e sim ampliar as possibilidades avaliativas por meio de diferentes instrumentos \textbf{que} contribuam na forma de expressar as habilidades desenvolvidas pelos estudantes, a partir da apropriação dos conhecimentos químicos ministrados em sala de aula.
Nessa perspectiva, o docente pode possibilitar alguns instrumentos de avaliação para os estudantes, tais como : dinâmicas por meio do lúdico, leitura e interpretação de textos de divulgação científica, produção escrita, leitura e interpretação de gráficos e tabelas, pesquisas, relatórios de atividades experimentais, apresentação de seminários, simulados on-line, uso de simuladores com situações contextualizadas, estratégias de argumentação como júri simulado, produção de vídeos e podcasts, infográficos, teatro, entre outros.
No ensino de Química, o foco da avaliação é compreender se o processo de letramento científico está sendo construído.
Portanto, independente da escolha dos instrumentos avaliativos, é importante \textbf{que} as questões problematizadoras estejam inseridas em contexto, e \textbf{que} seja avaliado o \textbf{raciocínio} do estudante durante todo o processo — não apenas o resultado final.
Dessa forma, é possível avaliar a leitura de mundo do estudante e se ele é capaz de utilizar o conhecimento escolar na resolução de problemas do cotidiano.
Avaliar o desenvolvimento do letramento científico na Educação Básica é pensar na formação integral do estudante objetivando \textbf{uma} reflexão sobre as relações das juventudes, sua participação política, visando a \textbf{uma} educação para a cidadania, com maior importância na \textbf{abordagem} do \textbf{sujeito}, numa perspectiva de proximidade entre o estudante e os novos padrões culturais.
Sendo assim, esta proposta não se limita a examinar os conhecimentos escolares apropriados, mas sim avaliar e desvelar os indivíduos, contribuindo para \textbf{que} o conhecimento adquirido avance no sentido de produzir informação \textbf{que} contribua com \textbf{uma} interação mais autônoma e crítica na sociedade.

% matched lemmas: abordagem, aplicabilidade, aproximar, autor, basear, categoria, central, conceitual, configurar, conseguir, contemporâneo, contextualização, correlação, crise, despertar, dialógico, direcionamento, dizer, encaminhamento, ensino-aprendizagem, erro, et, exigir, exposto, expressivo, focar, funcional, fundamento, historicamente, imprescindível, instigar, interdisciplinaridade, justo, lembrar, logo, memorização, mero, metodológico, moderno, método, necessitar, norteador, permear, pesquisador, posicionamento, preciso, problematização, prova, que, raciocínio, restringir, salientar, simples, sujeito, tarefa, teoria, teórico, totalidade, tópico, um, ver, viver, ótica
\end{textsample}
